\documentclass[11pt]{article}
\usepackage[margin=1in]{geometry}
\usepackage{amsmath,amssymb,amsthm,mathtools,microtype,hyperref}
\hypersetup{hidelinks}
\newtheorem{theorem}{Theorem}
\newtheorem{lemma}{Lemma}
\newcommand{\ddc}{dd^c}
\newcommand{\C}{\mathbb C}
\newcommand{\R}{\mathbb R}
\newcommand{\Z}{\mathbb Z}
\title{Prescribed logarithmic types on unibranch curves with torus normalization}
\author{}
\date{4 October 2026}
\begin{document}
\maketitle
\begin{abstract}
We give explicit sufficient conditions for equations in ordinary
singularity types on compact unibranch complex curves. Their normalization
is a complex torus and their Hermitian background is smooth in local
embeddings; its pullback can vanish at singular points. A periodic Green
function constructs every solution for normalization-side $L^p$ densities,
$p>1$, and a finite logarithmic pole type whose total weight is smaller
than the background mass. We prove finite mass for the entire genuine
ordinary type, uniqueness among initially rough solutions, holomorphic
disk descent, continuity of the residual, and independence of all
resolutions. An explicit projective curve has genus-one normalization,
one cusp, background mass six, a weight-one cusp pole, and a nonconstant equation
of mass five. The result is a dimension-one sufficient-condition answer
to P23; higher-dimensional nonclosed backgrounds remain outside its scope.
\end{abstract}

\section{Geometric hypotheses and the ordinary type}
Use $\ddc=(i/(2\pi))\partial\bar\partial$ throughout. In a flat
coordinate $z=x+iy$ this gives
\begin{equation}\label{eq:normalization}
 \ddc v=\frac{\Delta v}{4\pi}\,dx\,dy,
 \qquad \ddc\log|z|^2=\delta_0.
\end{equation}
The second identity follows from the divergence theorem on punctured
disks. All logarithmic coefficients below use squared logarithms.

Let $X$ be a connected compact reduced, locally irreducible complex
curve whose normalization
\[
 \pi:E=\C/\Lambda\longrightarrow X
\]
is finite and bijective, and is biholomorphic outside a finite set.
Let $h$ be any smooth Hermitian background in the local-embedding sense:
near every point it is the restriction of a smooth strictly positive
ambient $(1,1)$-form. Put $dA=dx\,dy$, $V=\int_EdA$, and
$\beta_0=dA/V$. Then
\[
 \beta=\pi^*h=b\beta_0,\qquad b\in C^\infty(E),\quad b\geq0,
 \qquad H=\int_E\beta>0.
\]
The pullback is positive over $X_{\rm reg}$ and can be degenerate elsewhere.
No positive pullback is assumed. Every smooth two-form on a curve is
closed, which is the dimension-one feature used in this paper.

Potentials are genuine canonical representatives. Specifically, a
potential on $X$ is tested on each holomorphic disk $f$: if $p_f$ is
a smooth local primitive of $f^*h$, then $u\circ f+p_f$ is
subharmonic or identically $-\infty$. A disk inside a pole is allowed
the latter alternative. The potential on $X$ is not identically
$-\infty$. Null-set changes of subharmonic values are excluded.

Choose distinct $q_1,\ldots,q_k\in E$ and $a_j>0$, with
\[
 A=\sum_ja_j<H,\qquad M=H-A.
\]
An empty list is allowed. We will construct a genuine reference $\psi$
whose pullback is $a_j\log|z-q_j|^2$ plus a smooth function near
$q_j$, and which has no other poles. Its ordinary type consists of
all genuine $h$-potentials $w$ satisfying
\begin{equation}\label{eq:ordinary-type}
 w\leq\psi+C,
 \qquad \psi\leq w+C
\end{equation}
for a finite constant $C$. These are extended-value inequalities; no
subtraction of two infinite values is used.

\begin{theorem}\label{thm:curve}
Under these hypotheses, let $0\leq f\in L^p(E,\beta_0)$ for some
finite $p>1$, with $F=\int_Ef\beta_0>0$, and put
$\nu=\pi_*(f\beta_0)$. This is a finite non-pluripolar measure.
Every genuine potential in the ordinary type \eqref{eq:ordinary-type}
has finite extended non-pluripolar mass $M$. There is a normalized
solution in that type of
\[
 \operatorname{MA}^{\rm ext}_h(u)=c\nu,\qquad
 \sup_Xu=0,\qquad c=\frac{H-A}{F}>0.
\]
All solutions in the entire genuine type, including initially rough
competitors with an unknown proportionality constant, differ by constants
and have this same $c$. The normalized solution is unique.

The potential is real-continuous away from its finite pole set, and
continuous into $\R\cup\{-\infty\}$ at the poles with the order
topology. Its residual relative to $\psi$, initially defined away
from the poles, extends to a continuous real function on all $X$.
The equation and type are independent of every resolution and common
refinement. In the bounded type $A=0$, the potential itself is
real-continuous everywhere.
\end{theorem}

The density hypothesis is measured against flat area on the normalization.
It is not a bounded-density assumption relative to the possibly
degenerating volume $h$ downstairs. The proof supplies descent and
rough-type uniqueness explicitly.

\section{The Green function and explicit construction}
There is a real mean-zero periodic distribution $G$ with
\begin{equation}\label{eq:green}
 \ddc G=\delta_0-\beta_0,
 \qquad G(z)=\log|z|^2+\text{smooth near }0,
\end{equation}
smooth away from $0$. To construct it, cut off $\log|z|^2$ in
a coordinate disk, with the cutoff equal to one near $0$, and call the
result $L$. Then $\ddc L=\delta_0+q\beta_0$, where $q$ is smooth
and $\int_Eq\beta_0=-1$, by pairing with the constant function.
Write the characters of $E$ as $e_\xi$, with
$-\Delta e_\xi=\lambda_\xi e_\xi$ and $\lambda_\xi>0$ when
$\xi\ne0$. The smooth mean-zero solution of
$\ddc S=-(q+1)\beta_0$ has coefficients
\[
 S_\xi=\frac{4\pi}{V}\frac{(q+1)_\xi}{\lambda_\xi}
 \quad(\xi\ne0),\qquad S_0=0.
\]
Integration by parts gives arbitrarily rapid decay of the numerator.
Thus the series and all derivatives converge uniformly. Set $G=L+S$
and subtract its mean. This proves \eqref{eq:green}, including its
local logarithmic form. In particular $G\in L^r(E,\beta_0)$ for
every finite $r$.

Convolution uses the group law and Haar probability $\beta_0$.
For $r\in L^p$ let $K=G*(r\beta_0)$. Differentiating against
smooth tests and applying Fubini gives
\begin{equation}\label{eq:convolution}
 \ddc K=\left(r-\int_Er\beta_0\right)\beta_0.
\end{equation}
For smooth $r$, Fourier coefficients show that $K$ is smooth. For
$p>1$, put $p'=p/(p-1)$. H\"older's inequality gives
\begin{align*}
 \|K\|_\infty&\leq\|G\|_{p'}\|r\|_p,\\
 |K(z+s)-K(z)|&\leq\|r\|_p
                      \|G(\cdot+s)-G\|_{p'}.
\end{align*}
Translations are continuous in $L^{p'}$: approximate $G$ in that norm
by smooth functions and use their uniform translation continuity.
Consequently $K$ is bounded and continuous. These estimates prove the
continuity needed here without a nonlinear existence or singular-space
regularity theorem.

Define
\[
 \Phi=G*((H-b)\beta_0),\qquad
 \Psi=\Phi+\sum_ja_jG(\cdot-q_j).
\]
Then $\Phi$ is smooth, and direct differentiation gives
\begin{equation}\label{eq:reference}
 \beta+\ddc\Psi=M\beta_0+\sum_ja_j\delta_{q_j}.
\end{equation}
Thus $\Psi$ is genuinely $\beta$-subharmonic. More explicitly, a
smooth primitive of $\beta$ exists in any coordinate disk by convolution
of its smoothly extended density with the logarithmic fundamental
solution. Adding it to $\Psi$ gives a subharmonic function with the
displayed positive Laplacian and the canonical logarithmic values.

Put $c=M/F$, and set
\begin{equation}\label{eq:solution}
 v=G*((cf-M)\beta_0),\qquad U=\Psi+v.
\end{equation}
The correction $v$ is bounded and continuous, and
\begin{equation}\label{eq:solution-current}
 \beta+\ddc U=cf\beta_0+\sum_ja_j\delta_{q_j}.
\end{equation}
It is positive. Local smoothing of its genuine potentials, together
with continuity off the poles and their explicit logarithmic values,
shows that $U$ is genuinely $\beta$-subharmonic. It belongs to the
ordinary type of $\Psi$. Its supremum $S$ is finite by compactness,
upper semicontinuity and nontriviality. We will descend $U-S$.

\section{Disk descent and continuity}
The finite bijective proper map $\pi$ is a homeomorphism between
compact Hausdorff spaces. Therefore the only possible descended
functions are
\[
 \psi=\Psi\circ\pi^{-1},\qquad u=(U-S)\circ\pi^{-1}.
\]
Fibers are singleton fibers, including at the poles.

Every holomorphic disk $f:D\to X$ has a holomorphic lift
$F=\pi^{-1}\circ f:D\to E$. For nonconstant $f$, the inverse
image of the finite singular set is discrete. The lift is holomorphic
off it. At an omitted point its continuous image stays in a bounded
coordinate neighborhood, so the usual removable-singularity theorem
for one-variable holomorphic functions extends the coordinate lift.
Constant disks have constant lifts; a nonconstant disk cannot lie
inside a finite singular set.

Locally choose a smooth primitive $p$ of $\beta$ on the image of $F$.
Then $p\circ F$ is a primitive of $f^*h$ and
\[
 u\circ f+p\circ F=(U-S+p)\circ F
\]
is subharmonic or identically $-\infty$. For completeness, composition
follows by decreasing smooth local regularizations of a subharmonic
function and the chain rule
$\Delta(a\circ F)=|F'|^2(\Delta a)\circ F$.
Any other smooth disk primitive differs by a harmonic function.
Thus the exact disk tests hold. A constant disk in a pole has the
allowed identically infinite restriction. A lifted nonconstant disk
meeting $q_j$ with order $m$ has pole coefficient $ma_j$, since
$\log|z(F(t))|^2=m\log|t|^2+\text{smooth}$.

The homeomorphism also proves the asserted continuity. Off the poles
$U$ is continuous, while its logarithmic form tends to $-\infty$
at each pole. The residual is the descended continuous real function
$v-S$ on all $X$; it is specified this way at the poles. This
description never forms $-\infty-(-\infty)$.

\section{Mass, full-type uniqueness and refinements}
Let $w$ be any genuine potential in the ordinary type and let
$W=w\circ\pi$. Applying the disk tests to normalization coordinate
disks proves genuine $\beta$-subharmonicity. The two type inequalities
give local boundedness on $E\setminus Q$, where $Q=\{q_j\}$,
and
\[
 W=a_j\log|z-q_j|^2+O(1)
\]
at each pole, with its canonical value $-\infty$. The positive
distribution $T_W=\beta+\ddc W$ is a positive Radon measure of total
mass $H$, by distributional Stokes on the compact torus.

Its atom at $q_j$ has precisely weight $a_j$. Indeed add a smooth
primitive of $\beta$ to $W$ and let $m(t)$ be its circular mean.
The divergence theorem, first for smooth regularizations and then in
distributions, gives for almost every $t$
\[
 m'(t)=\frac{2T_W(B_t)}{t}.
\]
If $\tau=T_W(\{q_j\})$, then $T_W(B_t)\to\tau$ as
$t\downarrow0$. Integrating gives
$m(t)/\log t^2\to\tau$. The logarithmic type gives the same
limit $a_j$, hence $\tau=a_j$. Coordinate changes do not alter the
weight because the corresponding logarithms differ by a smooth function.

The dimension-one truncation definition of the non-pluripolar product
now gives
\[
 \langle\beta+\ddc W\rangle=\mathbf1_{E\setminus Q}T_W.
\]
On every relatively compact set avoiding $Q$, a sufficiently low
truncation equals $W$; its linear product is $T_W$. The indicator
removes the poles. Therefore the entire type has mass $H-A=M$.
For the constructed solution \eqref{eq:solution-current}, removal
of the atoms leaves $cf\beta_0$, so
$\operatorname{MA}^{\rm ext}_h(u)=\pi_*(cf\beta_0)=c\nu$.

The measure $\nu$ is non-pluripolar. A nontrivial local
quasi-subharmonic function pulls back to a nontrivial locally
integrable quasi-subharmonic function on $E$. Its negative-infinite
set has flat area zero and hence zero $f\beta_0$-measure. The
homeomorphism transfers this property downstairs.

Suppose $w$ is any solution in the full ordinary type with
$\operatorname{MA}^{\rm ext}_h(w)=d\nu$, $d>0$. Its mass gives
$dF=M$, so $d=c$. The Borel homeomorphism and the fixed pole atoms
then imply
\[
 T_W=cf\beta_0+\sum_ja_j\delta_{q_j}=T_U.
\]
Thus $\ddc(W-U)=0$. All nonzero Fourier coefficients of this
$L^1$ distribution vanish by the flat Laplacian equation, so its
distribution is constant. Genuine local subharmonic representatives
are determined by their almost-everywhere classes: each value is the
limit of shrinking disk averages. Therefore $W=U+\text{constant}$
with canonical extended values, and the same equality holds downstairs.
This proves uniqueness among rough competitors as well.

Any resolution $\rho:Y\to X$ factors through $E$ by the same local
disk-lifting argument. The resulting proper map $r:Y\to E$ has
generic degree one. A holomorphic map between compact smooth curves
has finite fibers whose local multiplicities sum to its degree.
Degree one makes every multiplicity one and every fiber a singleton;
thus $r$ is biholomorphic. The same applies to every common refinement.
The type and measure are therefore unchanged on all resolutions.

The normalization has generic degree one even where its derivative
vanishes. There is no extra degree factor in $\nu=\pi_*(f\beta_0)$.
On the regular locus its density relative to $h$ is
$(f/b)\circ\pi^{-1}$; $b$ includes the derivative factor.
Both measures assign the finite singular set zero mass, and this density
formula is not used there. The degenerate pullback was compensated by
the explicit smooth correction $\Phi$, not declared strictly positive.

\section{An explicit projective cusp}
Take $E=\C/(\Z+i\Z)$, with $\beta_0=dx\,dy$, and define
\[
 x(z)=\frac1{z^2}+\sum_{\lambda\in\Z+i\Z,\ \lambda\ne0}
       \left(\frac1{(z-\lambda)^2}-\frac1{\lambda^2}\right),
 \qquad y=x'.
\]
The summands are $O(|\lambda|^{-3})$ on compacta away from the
lattice, so the series converges normally. Pairing opposite lattice
points gives evenness of $x$. The differentiated series is absolutely
normally convergent and reindexing proves periodicity of $y$.
Consequently $x(z+\omega)-x(z)$ is constant; evaluating at
$z=-\omega/2$, for $\omega=1,i$, and using evenness makes the
constant zero. Thus $x,y$ are meromorphic on $E$, with
\[
 x=z^{-2}+a_2z^2+a_4z^4+\cdots,
 \qquad y=-2z^{-3}+2a_2z+4a_4z^3+\cdots.
\]
There are no other poles. Laurent cancellation and the maximum
principle on $E$ give
\begin{equation}\label{eq:cubic}
 y^2=4x^3-g_2x-g_3,
 \qquad g_2=20a_2,\quad g_3=28a_4.
\end{equation}
The canceled elliptic function has no pole and constant Laurent
coefficient zero, proving the equality.

The argument principle on a translated fundamental parallelogram,
whose opposite boundary integrals cancel, shows that $x-c$ has
two zeros counted with multiplicity for every finite $c$.
Thus $x$ has degree two, and its fibers consist of $z,-z$ with
multiplicity. If $y(z)=0$, its local multiplicity is two and
$z=-z$ modulo the lattice; otherwise the two distinct points would
give at least four zeros. Also $x''(z)\ne0$ there. Since $y$ is odd,
the map $(x,y)$ is injective and immersive on $E\setminus\{0\}$:
use $dx$ when $y\ne0$ and $dy=x''dz$ when $y=0$.

Define
\[
 \pi(z)=[1:x(z):y(z):x(z)^2:x(z)^3]\in\mathbb P^4,
 \qquad \pi(0)=[0:0:0:0:1].
\]
In the last chart this holomorphic extension has coordinates
\[
 [s^3:s^2:t:s:1],\qquad s=1/x=z^2+O(z^6),
 \qquad t=y/x^3=-2z^3+O(z^7),
\]
and \eqref{eq:cubic} becomes
\[
 t^2=s^3(4-g_2s^2-g_3s^3).
\]
The factor in parentheses is a nonvanishing analytic unit near $0$.
Dividing $t$ by its analytic square root gives the ordinary cusp
$t'^2=s^3$, with local ring $\C\{\zeta^2,\zeta^3\}$ and
normalization $\C\{\zeta\}$, finite with module generators
$1,\zeta$. Indeed division by $t'^2-s^3$ leaves
$A(s)+t'B(s)$; the even and odd parts after substitution identify
the image ring. Its fraction field contains $\zeta$, which is integral
because $\zeta^2=s$. An integral meromorphic germ cannot have negative
order, since its highest power in a monic equation would have strictly
smallest order. Thus the integral closure is exactly $\C\{\zeta\}$.
Away from $0$ the map is injective and immersive, and
its first projective coordinate is nonzero. No such point maps to
$\pi(0)$.

The image is explicitly projective. In coordinates $Z_0,\ldots,Z_4$,
take the reduced projective locus
\begin{gather*}
 Z_1^2=Z_0Z_3,\qquad Z_1Z_3=Z_0Z_4,\qquad Z_3^2=Z_1Z_4,\\
 Z_2^2=4Z_0Z_4-g_2Z_0Z_1-g_3Z_0^2.
\end{gather*}
On $Z_0\ne0$ this is precisely the affine cubic graph with
coordinates $x^2,x^3$; the degree-two fibers of $x$ supply both
square-root values of $y$. On $Z_0=0$ the equations force
$Z_1=Z_3=Z_2=0$, leaving only $p$. In the last chart they give
exactly the cusp equation above. Hence this locus is $X=\pi(E)$,
smooth off $p$ and locally the stated reduced irreducible cusp there.
Thus $X$ is connected, unibranch, and has exactly one singular point
$p=\pi(0)$; $\pi$ is its finite bijective normalization.
It is not a projective cone. An irreducible reduced projective cone
of dimension one has a zero-dimensional irreducible reduced base,
hence is a projective line; our normalization has genus one and
our curve has a cusp. The cusp is nonnormal, so this example is
also outside log-terminal hypotheses.

Let $h$ be restricted Fubini--Study curvature with the convention
\eqref{eq:normalization}. Then $\beta=\pi^*h=b\beta_0$ is
smooth and positive off $0$, but
\[
 C_1|z|^2\leq b(z)\leq C_2|z|^2
\]
near $0$. All derivatives of the last-chart coordinates vanish
there, while $ds=2z\,dz+\cdots$ supplies the leading term.
The total background mass is $H=6$. Indeed on $E\setminus\{0\}$
the Fubini--Study potential
\[
 L=\log(1+|x|^2+|y|^2+|x|^4+|x|^6)
\]
satisfies $\ddc L=\beta$, while $L=-6\log|z|^2+\text{smooth}$
near $0$. Globally $\ddc L=\beta-6\delta_0$; pairing with one
proves the mass assertion. The full normalization has degree one,
although $x$ alone has degree two.

Choose one pole at $0$ with weight $a=1$. The reference construction
gives
\[
 \Phi=G*((6-b)\beta_0),\qquad \Psi=\Phi+G,
 \qquad \beta+\ddc\Psi=5\beta_0+\delta_0.
\]
For nonconstant data take
\[
 f(z)=1+\tfrac12\cos(2\pi\operatorname{Re}z),\qquad
 \nu=\pi_*(f\beta_0).
\]
Here $F=1$, $c=5$, and the continuous correction is explicitly
\[
 v(z)=-\frac5{2\pi}\cos(2\pi\operatorname{Re}z),
 \qquad \ddc v=(5f-5)\beta_0.
\]
The normalized descent of $\Psi+v$ satisfies
$\operatorname{MA}^{\rm ext}_h(u)=5\nu$ and has mass five.
The omitted unit mass is the pole atom. Every genuine potential in
this ordinary type has mass five, and all solutions are unique up
to constants. The potential tends to $-\infty$ at the cusp;
its residual is the continuous descended function $v-S$.
The density relative to $h$ behaves like $|z|^{-2}$ at the cusp
on the regular locus, so no bounded downstairs density is implicit.
With the empty pole list, the same curve instead gives $c=6$,
$v=-(3/\pi)\cos(2\pi\operatorname{Re}z)$, and a bounded-type
potential real-continuous on all $X$.

\section{Scope}
The explicit hypotheses are torus normalization, dimension one,
finite logarithmic ordinary type with $A<H$, and normalization-side
$L^p$ density for $p>1$. The proof treats arbitrary smooth local-embedding
backgrounds in that dimension, arbitrary genuine competitors in the
specified type, the degenerate pullback, and every resolution.
It proves a sufficient-condition result beyond projective cones,
without classifying arbitrary unibranch spaces or singularity types.
In particular it makes no higher-dimensional nonclosed existence or
regularity claim. The full P23 programme remains open.
\end{document}
